% Secao 4 - O deserto de distancia (descritivo).

\section{O deserto de distância}
\label{sec:distancia}

A medida que uma política de acesso baseada em mapa adota é direta: a distância
do município ao hospital ativo mais próximo. A Figura~\ref{fig:distance_map} a
traça para o Brasil. A mediana é de 16~km e, para metade dos municípios, o
hospital mais próximo está a menos de meia hora de estrada; só no quartil
superior a distância passa de 27~km, e apenas 5\% dos municípios estão a mais de
62~km.% src: distribuicao iso_km - p50 16, p75 27, p95 62, max 378
Lido isoladamente, o mapa de distância é tranquilizador: a maior parte do país
aparece em tons escuros, ``perto''.

A exceção visível é a Amazônia, onde a distância ao hospital mais próximo chega a
centenas de quilômetros. Mas é justamente ali que a medida é mais frágil. A
distância rodoviária pressupõe que o deslocamento se dá por estrada; no Norte,
ele se dá por rio, e a quilometragem sobre a malha viária superestima --- ou
simplesmente não representa --- a viagem real. O mapa de distância, portanto,
erra em dois sentidos que só ficarão visíveis quando confrontado com o fluxo:
exagera o isolamento onde o rio supre a estrada e o subestima onde o hospital
próximo não é o hospital usado.
